\section*{第 \ref{ch:g3:continents-and-oceans} 章——地球仪：大陆与海洋}

\begin{solution}{exo:g3:continents-and-oceans:1}
亚洲、非洲、北美洲、南美洲、\omterm{def:g3:continents-and-oceans:poles}{南极}洲、欧洲和
澳大利亚。最大的是亚洲；最小的是澳大利亚。
\end{solution}

\begin{solution}{exo:g3:continents-and-oceans:2}
太平洋、大西洋、印度洋、南大洋和北冰洋。
最大的是太平洋；最小的是北冰洋。
\end{solution}

\begin{solution}{exo:g3:continents-and-oceans:3}
\omterm{def:g3:continents-and-oceans:poles}{北极}位于\omterm{def:g3:continents-and-oceans:globe}{地球仪}的顶部；\omterm{def:g3:continents-and-oceans:poles}{南极}位于底部；
\omterm{def:g3:continents-and-oceans:equator}{赤道}绕在\omterm{def:g3:continents-and-oceans:globe}{地球仪}中间一圈。
\end{solution}

\begin{solution}{exo:g3:continents-and-oceans:4}
欧洲：\omterm{def:g3:continents-and-oceans:equator}{北半球}。澳大利亚：\omterm{def:g3:continents-and-oceans:equator}{南半球}。\omterm{def:g3:continents-and-oceans:poles}{南极}洲：\omterm{def:g3:continents-and-oceans:equator}{南半球}。北美洲：
\omterm{def:g3:continents-and-oceans:equator}{北半球}。
\end{solution}

\begin{solution}{exo:g3:continents-and-oceans:5}
北冰洋环绕\omterm{def:g3:continents-and-oceans:poles}{北极}。南大洋环绕
\omterm{def:g3:continents-and-oceans:poles}{南极}洲。
\end{solution}

\begin{solution}{exo:g3:continents-and-oceans:6}
从非洲到南美洲：大西洋。从澳大利亚向东到
南美洲：太平洋。从非洲到澳大利亚：印度
洋。
\end{solution}

\begin{solution}{exo:g3:continents-and-oceans:7}
山姆住在\omterm{def:g3:continents-and-oceans:equator}{赤道}附近，那里全年都很热。博住在
靠近极地的地方，那里全年都很冷。
\end{solution}

\begin{solution}{exo:g3:continents-and-oceans:8}
$10 - 7 = 3$：她的手指落在陆地上 3 次。这是我们可以
预期的：地球表面约 100 份中有 71 份是水，所以大约每 10 次中有 7 次
手指会停在海上。（只转 10 次时，她也可能
得到 6 或 8：每局的结果会有所不同。）
\end{solution}

\begin{solution}{exo:g3:continents-and-oceans:9}
为了画出\omterm{def:g3:continents-and-oceans:world-map}{世界地图}，沿穿过太平洋中部的一条线把\omterm{def:g3:continents-and-oceans:globe}{地球仪}剪开
并摊平。太平洋的一半落在\omterm{def:g2:drawing-to-map:map}{地图}的左边缘，另一半
落在右边缘。真实的地球上只有一个太平洋：\omterm{def:g2:drawing-to-map:map}{地图}的
两条边缘在\omterm{def:g3:continents-and-oceans:globe}{地球仪}上是同一条
线。
\end{solution}

\begin{solution}{exo:g3:continents-and-oceans:10}
汤姆说错了。按从大到小排列，各\omterm{def:g3:continents-and-oceans:continent}{大陆}依次为亚洲、
非洲、北美洲、南美洲、\omterm{def:g3:continents-and-oceans:poles}{南极}洲、欧洲和澳大利亚：
\omterm{def:g3:continents-and-oceans:poles}{南极}洲只排第五，远小于亚洲。\omterm{def:g3:continents-and-oceans:world-map}{世界地图}
之所以误导了他，是因为平面\omterm{def:g2:drawing-to-map:map}{地图}会把靠近极地的陆地拉得很长。
\omterm{def:g3:continents-and-oceans:poles}{南极}洲环绕\omterm{def:g3:continents-and-oceans:poles}{南极}，在\omterm{def:g2:drawing-to-map:map}{地图}上被拉成横跨整个底部的一条长带。
在\omterm{def:g3:continents-and-oceans:globe}{地球仪}上，它是一块圆形的
陆地，远小于亚洲。
\end{solution}
